\documentclass[11pt]{article}
\usepackage[margin=1in]{geometry}
\usepackage{amsmath,amssymb,amsfonts,mathtools}
\usepackage{array,booktabs,enumitem,graphicx,xcolor,hyperref}
\usepackage[T1]{fontenc}
\usepackage[utf8]{inputenc}
\title{Batalin-Vilkovisky Lagrangian Quantisation}
\author{Source-backed arXiv sample}
\date{}
\begin{document}
\maketitle
\section{Extracted Lists}

The list environments below are extracted from arXiv LaTeX source and wrapped for pdfLaTeX validation.

\begin{itemize}
\item $S_{cl}(\phi ) = S(\Phi ,0)$~: the classical
action is recovered when all antifields are set to zero.
\item The extended action should satisfy the {\bf master equation}, which
at the classical level reads $(S,S)=0$.
\item $S$ is a proper solution, which means that $S_{\alpha \beta }\equiv
\dl_\alpha \dr_\beta S$ is
a matrix of rank $N$ on the stationary surface where $N$ is the number of
fields $\Phi ^A$.
\end{itemize}

\begin{itemize}
\item Point transformations are the easiest ones. This are just
redefinitions between the fields $\Phi '^A=f^A(\Phi )$. They are obtained
by $F=\Phi '^*_A f^A(\Phi )$ which thus determines the corresponding
transformations of the antifields. The latter replace the calculations of
the variations of the new variables.
\item Adding the BRST transformation of a function $s\Psi(\Phi )$ to the
action is obtained by a canonical transformation with $F=\Phi '^*_A\Phi ^A+
\Psi (\Phi )$. The latter gives
\begin{equation}
\Phi'^A = \Phi^A \ ;\qquad
\Phi ^*_A =\Phi'^*_A + \partial _A \Psi (\Phi ).
\label{gfermion}\end{equation}
\item Redefine the symmetries by adding equation of motions
('trivial symmetries'). This is obtained by
\begin{equation}
F=\Phi '^*_A\Phi ^A+ \Phi '^*_A\Phi '^*_B h^{AB}(\Phi )
\end{equation}
(the first term is the identity transformation).
\item Elimination of auxiliary fields can be done by canonical
transformations
(see appendix B of \cite{disp}). This procedure will then also give
the 'compensating transformations'. \end{itemize}
\end{document}
